% ============================================================
%  IKK III — Matrix Mechanics / Wave Function Bridge
%  Projected Wave Form, Matrix-State Representation,
%  and Recoverable Identity
%  VSIM Framework  —  Version 2.0
% ============================================================

\documentclass[12pt, a4paper]{article}

\usepackage{amsmath, amssymb, amsthm}
\usepackage{mathrsfs}
\usepackage{mathtools}
\usepackage{geometry}
\usepackage{hyperref}
\usepackage{booktabs}
\usepackage{array}
\usepackage{xcolor}
\usepackage{enumitem}
\usepackage{microtype}
\usepackage{lmodern}
\usepackage[T1]{fontenc}
\usepackage{bm}

\geometry{top=2.8cm, bottom=2.8cm, left=3.2cm, right=3.2cm}

\hypersetup{colorlinks=true, linkcolor=black, citecolor=black, urlcolor=black}

% ------------------------------------------------------------------
%  Theorem environments
% ------------------------------------------------------------------
\theoremstyle{definition}
\newtheorem{defn}{Definition}[section]

\theoremstyle{plain}
\newtheorem{prop}{Proposition}[section]
\newtheorem{thm}{Theorem}[section]
\newtheorem{construction}{Construction}[section]
\newtheorem{claim}{Claim}[section]

\theoremstyle{remark}
\newtheorem{remark}{Remark}[section]
\newtheorem{note}{Note}[section]
\newtheorem{warning}{Warning}[section]

\newenvironment{epistemic}
  {\medskip\noindent\textit{Epistemic status.}\enspace}{\medskip}

\newenvironment{evidenceitem}[1]
  {\medskip\noindent\textbf{#1.}\enspace}{\medskip}

% ------------------------------------------------------------------
%  Notation macros
% ------------------------------------------------------------------

%%  Identity objects
\newcommand{\Itrue}{\mathfrak{I}_{\mathrm{true}}}
\newcommand{\Itop}{\mathfrak{I}_{\mathrm{top}}}
\newcommand{\Ia}{\mathfrak{I}_a}
\newcommand{\Ihat}{\widehat{\mathfrak{I}}_a}
\newcommand{\II}{\mathfrak{I}}

%%  Spaces
\newcommand{\Ha}{\mathfrak{H}_a}        % source identity space
\newcommand{\Hb}{\mathcal{H}_b}         % target Hilbert space
\newcommand{\Hahat}{\widehat{\mathfrak{H}}_a}  % reconstruction codomain
\newcommand{\Kab}{\mathfrak{K}_{ab}}    % coupled cross-scale space
\newcommand{\Osp}{\mathcal{O}_b}        % target observable space
\newcommand{\ProbR}{\mathcal{P}(\mathbb{R}^3)} % probability measures over R^3

%%  Scale labels
\newcommand{\Sa}{S_a}
\newcommand{\Sb}{S_b}

%%  Projection and related operators
\newcommand{\Pab}{\mathcal{P}_{a\to b}}
\newcommand{\Wab}{\mathcal{W}_{a\to b}}   % wave-state recovery map
\newcommand{\Qba}{\mathcal{Q}_{b\to a}}   % reconstruction operator
\newcommand{\Cab}{\mathcal{C}_{ab}}        % cross-scale coupling
\newcommand{\Rb}{\mathcal{R}_b}            % resolution map
\newcommand{\Mb}{\mathcal{M}_b}            % general measurement map
\newcommand{\MXb}{\mathcal{M}_{X,b}}      % position measurement map
\newcommand{\Oba}{O_b^{(a)}}

%%  Hilbert space and wave objects
\newcommand{\psiab}{|\psi_b^{(a)}\rangle}
\newcommand{\psibra}{\langle\psi_b^{(a)}|}
\newcommand{\psix}{\psi_b^{(a)}(x)}
\newcommand{\psirec}{\psi_{b,\mathrm{rec}}^{(a)}}
\newcommand{\psiQM}{\psi_{\mathrm{QM}}}
\newcommand{\Aab}{A_b^{(a)}}
\newcommand{\phiab}{\phi_b^{(a)}}

%%  Position and spectral measure
\newcommand{\Xhat}{\hat{\mathbf{X}}}
\newcommand{\EX}{E_X}
\newcommand{\muX}{\mu_{X,b}^{(a)}}
\newcommand{\Xexp}{\langle\mathbf{X}\rangle_b^{(a)}}

%%  Density
\newcommand{\rhoab}{\rho_b^{(a)}}
\newcommand{\rhoQM}{\rho_{\mathrm{QM}}}
\newcommand{\rhorec}{\rho_{b,\mathrm{rec}}^{(a)}}

%%  Identity state vector and evolution
\newcommand{\sa}{\mathbf{s}_a}
\newcommand{\Ua}{\hat{U}_a}
\newcommand{\Hhat}{\hat{H}}
\newcommand{\Uhat}{\hat{U}}

%%  Distinguishability functionals
\newcommand{\Dista}{\mathfrak{D}_a}
\newcommand{\Distb}{\mathfrak{D}_b}
\newcommand{\Drec}{\mathfrak{D}_{\mathrm{rec},b}^{(a)}}

%%  Projection loss quantities
\newcommand{\chiab}{\chi_{ab}}
\newcommand{\etaab}{\eta_{ab}}
\newcommand{\etamin}{\eta_{\min}}
\newcommand{\epsI}{\varepsilon_{\mathfrak{I},ab}}
\newcommand{\dHa}{d_{\Ha}}
\newcommand{\Lpsi}{\mathcal{L}_\psi}
\newcommand{\Lrho}{\mathcal{L}_\rho}

%%  Bilinear kernel objects
\newcommand{\Rbvec}{R_b(r)}
\newcommand{\eavec}{e_a(t)}
\newcommand{\Mab}{M_{ab}}
\newcommand{\Bbase}{B(r)}
\newcommand{\Dmod}{D(r)}

%%  Scale coordinate
\newcommand{\uscl}{u}
\newcommand{\Duab}{\Delta u_{ab}}

% ------------------------------------------------------------------
\begin{document}
% ------------------------------------------------------------------

\begin{titlepage}
  \centering
  \vspace*{2cm}
  {\Large\scshape VSIM Framework\par}
  \vspace{0.4cm}
  {\large IKK Series\,—\,Inverse Kaluza-Klein / Scale-Projection\par}
  \vspace{1.2cm}
  {\Huge\bfseries IKK III\par}
  \vspace{0.5cm}
  {\LARGE Matrix Mechanics / Wave Function Bridge\par}
  \vspace{0.4cm}
  {\large\itshape Projected Wave Form, Matrix-State Representation,\\
  and Recoverable Identity\par}
  \vspace{1.8cm}
  {\large Version 2.0\par}
  \vfill
  {\small
    \textit{Prerequisite documents: Document~0, IKK~I, IKK~II.}\par
    \vspace{0.4cm}
    \textit{Version 2.0 introduces typed identity and Hilbert spaces,
    the wave-state recovery map $\Wab$, spectral position measures,
    a reconstruction operator, and rigorous efficiency functionals.
    The core interpretation is unchanged; the mathematical infrastructure
    is upgraded throughout.}\par
    \vspace{0.4cm}
    \textit{Epistemic status: proposed reinterpretation.
    The standard quantum mechanical formalism is accepted as a
    correct predictive framework at every scale where it has been tested.}
  }
\end{titlepage}

\tableofcontents
\newpage

% ==================================================================
\section*{Central Claim}
\addcontentsline{toc}{section}{Central Claim}
% ==================================================================

In IKK~III, the wavefunction is not postulated as the primitive identity
object.
It is modelled as the target-scale Hilbert-space recovery of source
identity content:

\begin{center}
\fbox{\parbox{0.82\textwidth}{\centering\vspace{0.3cm}
  \[
    |\psi_b^{(a)}\rangle
    = (\Rb\circ\Cab)(\Ia)
    = \Wab(\Ia)
  \]
  \vspace{0.1cm}
  \textit{The wavefunction is what identity looks like after it has
  been carried through a typed cross-scale coupling and a typed
  resolution map into a target Hilbert space.}
\vspace{0.3cm}}}
\end{center}

\medskip
The primitive order is:
\[
  \Ia \in \Ha
  \;\xrightarrow{\;\Cab\;}\;
  \Kab
  \;\xrightarrow{\;\Rb\;}\;
  \Hb
  \;\ni\; \psiab
\]
Each arrow is a typed map between named spaces.
$\Ha$ is the source identity space.
$\Hb$ is the target Hilbert space.
They are not identified until the map $\Wab = \Rb\circ\Cab$ has
been applied.

\medskip
Position becomes a spectral measure:
\[
  \muX(\Omega)
  = \langle \Wab(\Ia),\,\EX(\Omega)\,\Wab(\Ia)\rangle_{\Hb}
\]
Reconstruction carries an error:
\[
  \epsI = \dHa(\Ia,\,\Qba\Pab(\Ia))
\]
Both are definitions, not claims. Evaluating them is the work of the EVM series.


% ==================================================================
\section{Purpose}
\label{sec:purpose}
% ==================================================================

This document explains how matrix mechanics and wave mechanics can be
treated as two different representational layers of the same underlying
identity structure, and how that structure is made precise through
typed maps between named spaces.

The standard state evolution:
\[
  i\hbar\frac{\partial}{\partial t}\psi(x,t) = \Hhat\,\psi(x,t),
  \qquad
  |\psi(t)\rangle = \Uhat(t)|\psi(0)\rangle
\]
is accepted as correct within its domain.
The dispute is interpretive.

\begin{align*}
  \text{matrix mechanics} &= \text{structured identity-state update} \\[4pt]
  \text{wave mechanics}   &= \text{projected observable recovery}
\end{align*}

The bridge connecting these is:
\[
  \Ia(t)
  \;\longrightarrow\; \sa(t)
  \;\longrightarrow\; \Ua\,\sa(t)
  \;\longrightarrow\; \Wab[\Ia(t)]
  \;\longrightarrow\; \psiab
\]

Version 2.0 upgrades the framework in ten respects.
The projection is now a typed map between named spaces, not a
suggestive arrow.
The Hilbert space $\Hb$ at the target scale is explicitly separated
from the identity space $\Ha$ at the source.
Position is a spectral measure induced by projection, not a
recovered coordinate.
Reconstruction is a separate map with a quantified error.
Efficiency $\etaab$ is a functional on $\Ha$, not a scalar in $[0,1]$.
Compatibility $\chiab$ is a continuous gate, with the Boolean version
as a limiting case.

\begin{remark}
The logarithmic scale coordinate $\uscl = \ln(\ell/\ell_0)$ is used
throughout as the canonical scale variable.
The rhetorical fifth axis $w$ introduced in Document~0 remains valid
as an introductory device; it is not used as a formal variable here.
\end{remark}


% ==================================================================
\section{The Historical Split and What It Exposes}
\label{sec:history}
% ==================================================================

Matrix mechanics gives transition structure and operator algebras:
\[
  |\psi(t+\Delta t)\rangle = \Uhat(\Delta t)\,|\psi(t)\rangle.
\]
Wave mechanics gives the spatially resolved amplitude form:
\[
  i\hbar\,\partial_t\psi = \Hhat\,\psi.
\]
Their predictive equivalence is established.
IKK~III asks the prior question:

\begin{center}
\textit{What are these equivalent descriptions descriptions of?}
\end{center}

\begin{prop}[Representational fracture]
\label{prop:fracture}
Predictive equivalence does not imply ontological identity:
\[
  \text{equivalent predictions} \;\neq\; \text{identical ontology.}
\]
Matrix mechanics behaves like a transition algebra over identity
state vectors:
\[
  \sa(t+\Delta t) = \Ua(\Delta t)\,\sa(t).
\]
Wave mechanics behaves like a resolved observable field at target
scale $\Sb$:
\[
  |\psi_b^{(a)}(t)\rangle = \Wab(\Ia(t)).
\]
The fracture between these two behaviours is meaningful.
It reflects the distinction between internal identity-state update
and observable recovery through projection.
\end{prop}


% ==================================================================
\section{Identity Spaces and the Type Separation}
\label{sec:spaces}
% ==================================================================

\subsection{Source identity space and target Hilbert space}

\begin{defn}[Source identity space]
\label{def:Ha}
The \emph{source identity space} $\Ha$ is the space in which source
identity content $\Ia$ resides.
It is the domain of the projection operator.
No inner product or Hilbert structure is imposed on $\Ha$ at this stage.
\end{defn}

\begin{defn}[Target Hilbert space]
\label{def:Hb}
The \emph{target Hilbert space} $\Hb$ is the standard quantum
mechanical Hilbert space at target scale $\Sb$.
It carries the inner product $\langle\cdot,\cdot\rangle_{\Hb}$,
observables as self-adjoint operators, and the spectral calculus.
\end{defn}

\begin{warning}
$\Ha$ and $\Hb$ must not be identified until the map $\Wab$ has been
applied.
Identifying them collapses the framework back into standard quantum
mechanics while pretending it has added structure.
The separation is the point.
\end{warning}

\subsection{The three-level identity separation}

Three objects are distinct and must not be conflated:

\begin{defn}[Three-level identity separation]
\label{def:threelevel}
\begin{enumerate}[label=(\roman*)]
  \item $\Itrue \in \Ha$ — the full primitive identity content.
  \item $\Itop$ — the highest packed or top-layer identity.
  \item $\Oba \in \Osp$ — the observable at target scale $\Sb$
    recovered from source $\Sa$.
\end{enumerate}
\[
  \Itrue \;\neq\; \Itop \;\neq\; \Oba
\]
\end{defn}

Classical quantum mechanics treats $|\psi\rangle$ as the complete state.
IKK~III requires:
\[
  |\psi\rangle \not\equiv \Itrue.
\]
The correct identification is:
\[
  |\psi_b^{(a)}\rangle \;\sim\; \Oba,
  \qquad
  \Oba = \Pab(\Ia).
\]

\subsection{Typed projection operator}

\begin{defn}[Typed projection operator]
\label{def:typedproj}
The \emph{projection operator} is a typed map:
\[
  \boxed{\Pab : \Ha \to \Osp}
\]
with:
\[
  \Ia \in \Ha,
  \qquad
  \Oba \in \Osp,
  \qquad
  \Oba = \Pab(\Ia).
\]
\end{defn}

\begin{remark}
Projection is no longer a decorative arrow.
It is a map between two named spaces with declared domain and codomain.
This makes the statement $\Oba = \Pab(\Ia)$ a typed functional equation,
not a notational convenience.
\end{remark}


% ==================================================================
\section{The Wave-State Recovery Map}
\label{sec:wavemap}
% ==================================================================

\subsection{Definition}

\begin{defn}[Wave-state recovery map]
\label{def:Wab}
The \emph{wave-state recovery map} is:
\[
  \boxed{\Wab : \Ha \to \Hb}
\]
It carries source identity $\Ia \in \Ha$ into a quantum state
$\psiab \in \Hb$ at target scale $\Sb$:
\[
  \boxed{|\psi_b^{(a)}\rangle = \Wab(\Ia)}
\]
\end{defn}

\begin{defn}[Position-basis wavefunction]
\label{def:psix}
Given $|\psi_b^{(a)}\rangle = \Wab(\Ia) \in \Hb$, the
position-basis wavefunction is the standard inner product:
\[
  \boxed{\psix = \langle x\,|\,\psi_b^{(a)}\rangle}
\]
The IKK composition is therefore:
\[
  \boxed{\psix = \langle x\,|\,\Wab(\Ia(t))\rangle}
\]
\end{defn}

\begin{remark}
This separates two operations that should not be merged.
$\Wab$ carries identity content from $\Ha$ into $\Hb$.
The inner product $\langle x\,|\cdot\rangle$ then evaluates the
resulting state at a position.
These are distinct steps.
Writing $\psi \sim \mathcal{P}[\mathfrak{I}]$ merged them silently.
\end{remark}

\subsection{Relation to the projection operator}

The wave-state recovery map is the composition of the typed
cross-scale chain (formalised in Section~\ref{sec:typedchain}):
\[
  \boxed{\Wab = \Rb \circ \Cab : \Ha \to \Hb}
\]
so that:
\[
  |\psi_b^{(a)}\rangle
  = (\Rb \circ \Cab)(\Ia)
\]
and the position wavefunction is:
\[
  \psix = \langle x\,|\,(\Rb \circ \Cab)(\Ia)\rangle.
\]


% ==================================================================
\section{Typed Projection Chain}
\label{sec:typedchain}
% ==================================================================

\subsection{Three typed maps}

\begin{defn}[Cross-scale coupling]
\label{def:Cab}
\[
  \boxed{\Cab : \Ha \to \Kab}
\]
$\Kab$ is the \emph{coupled cross-scale information space}.
$\Cab$ determines which identity content from $\Ha$ is physically
accessible in the cross-scale interaction with $\Sb$.
This is the access condition, not yet an observable.
\end{defn}

\begin{defn}[Resolution map]
\label{def:Rb}
\[
  \boxed{\Rb : \Kab \to \Hb}
\]
$\Rb$ maps accessible identity content into the target Hilbert space.
It encodes what the target scale can represent, independently of any
measurement apparatus.
\end{defn}

\begin{defn}[Position measurement map]
\label{def:MXb}
\[
  \boxed{\MXb : \Hb \to \ProbR}
\]
where $\ProbR$ denotes the space of probability measures over
position space $\mathbb{R}^3$.
$\MXb$ extracts position probability information from a state
in $\Hb$ via the spectral measure of the position operator.
\end{defn}

\begin{warning}
$\Rb \neq \MXb$.
Resolution is the representational capacity of the target domain.
Measurement is what an apparatus, basis, or spectral decomposition
extracts from that capacity.
These are distinct operations.
\end{warning}

\subsection{Composition and typing}

The full position projection is the typed composition:
\[
  \boxed{\MXb \circ \Rb \circ \Cab : \Ha \to \ProbR}
\]
with staged information loss at each step.
Let $\mathfrak{D}$ denote recoverable distinguishability content.
Then:
\[
  \Dista(\Ia)
  \;\geq\; \mathfrak{D}_{\Kab}(\Cab(\Ia))
  \;\geq\; \mathfrak{D}_{\Hb}(\Rb(\Cab(\Ia)))
  \;\geq\; \mathfrak{D}_{\ProbR}(\MXb(\cdots))
\]
Each inequality may be strict.

\subsection{Uncertainty as accumulated projection loss}

The Heisenberg uncertainty relation:
\[
  \Delta x\,\Delta p \;\geq\; \frac{\hbar}{2}
\]
is reframed as a bound on jointly recoverable projected observables
at $\Sb$:
\[
  \Delta x_b^{(a)}\,\Delta p_b^{(a)} \;\geq\; \frac{\hbar}{2}.
\]
This locates uncertainty at the recovered observable layer.
It does not assert that $\Itrue$ is classically deterministic.

\begin{remark}
This framing prevents the standard failure mode: asserting
hidden-variable determinism at the identity level without providing
a mechanism that is experimentally distinguishable from standard
quantum mechanics.
Uncertainty belongs to the projection layer unless the primitive
layer is independently and specifically constrained.
\end{remark}


% ==================================================================
\section{Position as a Spectral Measure}
\label{sec:position}
% ==================================================================

\subsection{Position operator and spectral measure}

In standard quantum mechanics, position is a self-adjoint operator
$\Xhat$ on $\Hb$ with spectral measure $\EX$.
For a Borel set $\Omega \subseteq \mathbb{R}^3$, $\EX(\Omega)$ is
the projection onto states localised in $\Omega$.
The probability of finding the particle in $\Omega$ given state
$|\phi\rangle \in \Hb$ is:
\[
  \Pr(\Omega) = \langle \phi \,|\, \EX(\Omega) \,|\, \phi\rangle.
\]

\subsection{Projected position measure}

\begin{prop}[Projected position distribution]
\label{prop:positionmeasure}
Let $\Ia \in \Ha$ be a source identity state and let
$\Wab : \Ha \to \Hb$ be the wave-state recovery map of
Definition~\ref{def:Wab}.
Let $\EX$ be the spectral measure of the position operator
$\Xhat$ on $\Hb$.

The \emph{target-scale position distribution} of $\Ia$ is the
probability measure $\muX$ on $\mathbb{R}^3$ defined by:
\[
  \boxed{
    \muX(\Omega)
    = \bigl\langle \Wab(\Ia),\;
      \EX(\Omega)\,\Wab(\Ia)
      \bigr\rangle_{\Hb}
  }
\]
for every Borel set $\Omega \subseteq \mathbb{R}^3$.

If $\Wab(\Ia) = \psiab$ and $\psix = \langle x\,|\psi_b^{(a)}\rangle$,
this reduces to the Born rule:
\[
  \boxed{
    \muX(\Omega) = \int_{\Omega} |\psix|^2\,d^3x.
  }
\]
\end{prop}

\begin{remark}
Position is not a recovered coordinate $x_b^{(a)}$ assigned to
identity $\Ia$.
It is the probability measure $\muX$ induced on $\mathbb{R}^3$
by projecting $\Ia$ into $\Hb$ and applying the spectral measure
of $\Xhat$.
The Born rule emerges as the special case where $\Hb$ is a standard
$L^2$ Hilbert space and $\Wab(\Ia)$ is a normalised vector.
\end{remark}

\subsection{Explicit composition form}

Substituting $\Wab = \Rb \circ \Cab$:
\[
  \boxed{
    \muX(\Omega)
    = \Bigl\langle
        (\Rb\circ\Cab)(\Ia)
        \;\Big|\;
        \EX(\Omega)
        \;\Big|\;
        (\Rb\circ\Cab)(\Ia)
      \Bigr\rangle_{\Hb}
  }
\]
This is the formal bridge.
The position observable at $\Sb$ is entirely determined by the
typed composition of coupling, resolution, and spectral measure.


% ==================================================================
\section{Expected Position}
\label{sec:exppos}
% ==================================================================

Once the position measure $\muX$ is defined, the expected position
is the standard Lebesgue integral:
\[
  \boxed{
    \Xexp = \int_{\mathbb{R}^3} \mathbf{x}\;d\muX(\mathbf{x})
  }
\]
In Hilbert-space form:
\[
  \boxed{
    \Xexp = \langle\psi_b^{(a)}\,|\,\Xhat\,|\,\psi_b^{(a)}\rangle
  }
\]
IKK substitution:
\[
  \boxed{
    \Xexp
    = \Bigl\langle \Wab(\Ia) \;\Big|\; \Xhat \;\Big|\; \Wab(\Ia) \Bigr\rangle
  }
\]

\begin{remark}
Position is not the primitive coordinate of identity.
Position is the expectation structure of the target-scale position
observable induced by the projection of identity into $\Hb$.
The distinction is not terminological.
Treating position as a primitive coordinate of $\Ia$ would require
$\Ia$ to already live in $\mathbb{R}^3$, collapsing the framework
back into classical mechanics before the projection has done
any work.
\end{remark}


% ==================================================================
\section{Wavefunction Shape: Amplitude and Phase}
\label{sec:waveshape}
% ==================================================================

The wavefunction has wave-like form because the observer accesses
$\Wab(\Ia) \in \Hb$, not $\Ia \in \Ha$ directly.

\begin{defn}[Projected amplitude and phase]
\label{def:ampphase}
In the amplitude-phase decomposition
$|\psi_b^{(a)}(x,t)| = \Aab(x,t)\,e^{i\phiab(x,t)}$:
\begin{itemize}
  \item $\Aab(x,t)$ is the recovered amplitude of projected identity
    content at $\Sb$.
  \item $\phiab(x,t)$ is the phase relation induced by cross-scale
    coupling $\Cab$, resolution geometry $\Rb$, and the observation
    basis.
\end{itemize}
\[
  \Aab,\;\phiab \;\in\; \Hb
  \qquad\text{but}\qquad
  \Aab,\;\phiab \;\notin\; \Ha.
\]
\end{defn}

\begin{claim}[Wave is the shape of projection]
\label{claim:waveshape}
Wave-like behavior appears in the observable layer because
$\Wab$ maps structured identity content in $\Ha$ into
a spatially resolved amplitude-phase form in $\Hb$.
It does not follow that the primitive identity $\Ia$ is wave-substance.
The wave is the shape of the projection, not the shape of the identity.
\end{claim}

The probability density at $\Sb$:
\[
  \rhoab(x,t) = |\psix|^2
  = |\langle x\,|\,\Wab(\Ia(t))\rangle|^2
\]
recovers the Born rule as a consequence of the spectral measure
of Definition~\ref{prop:positionmeasure}, not as a postulate.


% ==================================================================
\section{Bilinear Kernel: Worked Projection Example}
\label{sec:kernel}
% ==================================================================

\subsection{Motivation}

Sections~\ref{sec:spaces}--\ref{sec:waveshape} establish typed maps
and spectral measures.
This section provides a tractable computational model for how
$\Wab$ acts in practice.

\subsection{Density recovery formula}

\begin{defn}[Target-scale basis vector]
\label{def:Rbvec}
$\Rbvec = [R_{b,1}^2(r),\ldots,R_{b,n}^2(r)]^T \in \mathbb{R}^n$
is the vector of squared radial basis functions at scale $\Sb$.
\end{defn}

\begin{defn}[Source identity channel vector]
\label{def:eavec}
$\eavec = [1, e_{a,1}(t),\ldots,e_{a,m}(t)]^T \in \mathbb{R}^{m+1}$
is the accessible state-channel vector of $\Ia$.
The leading entry $1$ anchors the static background component.
\end{defn}

\begin{defn}[Channel mixing matrix]
\label{def:Mab}
$\Mab \in \mathbb{R}^{n\times(m+1)}$ encodes the projection weight
from each source channel to each target basis component.
It is the matrix representation of $\Wab$ in the chosen bases.
\end{defn}

The density recovery formula is:
\[
  \boxed{\rhoab(r,t) = \Rbvec^T\,\Mab\,\eavec}
\]
This is the density form of $\Oba = \Pab(\Ia)$ under the bilinear
approximation, consistent with the Born rule:
\[
  |\langle x\,|\,\Wab(\Ia)\rangle|^2 \;\approx\; \Rbvec^T\,\Mab\,\eavec.
\]

\begin{remark}
$\Mab$ is determined by the identity structure and the choice of
basis, not fitted to observed data after the fact.
Fitting $\Mab$ post hoc to minimise recovery error produces a
curve-fitting exercise, not a projection model.
\end{remark}

\subsection{Cacheable separability}

For the two-basis case ($n = 2$):
\[
  \Bbase = \tfrac{1}{2}(R_1^2(r)+R_2^2(r)),
  \qquad
  \Dmod  = \tfrac{1}{2}(R_1^2(r)-R_2^2(r)).
\]
With a single dynamical channel $e_z(t)$:
\[
  \rho(r,t) = \Bbase + e_z(t)\,\Dmod,
  \qquad
  \nabla\rho(r,t) = \nabla\Bbase + e_z(t)\,\nabla\Dmod.
\]

\begin{prop}[Cacheable spatial structure]
\label{prop:cache}
$\Bbase$ and $\Dmod$ are time-independent and may be precomputed.
State evolution propagates only the scalar channel $e_z(t)$.
The full density and its gradient are recovered at each time step
without recomputing the basis structure.
\end{prop}

This is where the projection bridge becomes computationally operational.
Internal identity update propagates a small vector;
spatial structure is precomputed once.


% ==================================================================
\section{Projection Loss and Efficiency}
\label{sec:loss}
% ==================================================================

\subsection{Distinguishability as a functional}

\begin{defn}[Distinguishability functionals]
\label{def:dist}
Let:
\[
  \boxed{\Dista : \Ha \to \mathbb{R}_{\geq 0}}
\]
be the \emph{source distinguishability functional}, measuring the
identity information content of $\Ia \in \Ha$.
Let:
\[
  \boxed{\Distb : \Osp \to \mathbb{R}_{\geq 0}}
\]
be the \emph{observable distinguishability functional} at $\Sb$.
\end{defn}

\subsection{Projection efficiency}

\begin{defn}[Projection efficiency functional]
\label{def:eta}
The \emph{projection efficiency} of $\Ia$ under $\Pab$ is:
\[
  \boxed{
    \etaab(\Ia)
    = \frac{\Distb(\Pab(\Ia))}{\Dista(\Ia)}
  }
\]
with $0 \leq \etaab(\Ia) \leq 1$ when $\Distb$ is normalised
consistently with $\Dista$.
\end{defn}

\begin{remark}
$\etaab$ is now a functional on $\Ha$, not a scalar in $[0,1]$.
It depends on the specific identity state $\Ia$, not only on the
pair of scales.
Two identity states at the same source scale may project with
different efficiency depending on their structure relative to $\Cab$.
\end{remark}

\subsection{Scale dependence}

Using the logarithmic scale coordinate $\uscl = \ln(\ell/\ell_0)$,
the scale separation is $\Duab = |\ln(\ell_b/\ell_a)|$.
The framework should predict:
\[
  \etaab(\Ia) = f(\Duab,\,\Ia)
\]
for some function $f$ determined by the projection structure,
not fitted freely per case.

\subsection{Recovery loss functionals}

\begin{defn}[Recovery loss]
\label{def:loss}
\[
  \Lpsi = \int \bigl|\psiQM(x,t) - \psirec(x,t)\bigr|^2\,dx,
  \qquad
  \Lrho = \int \bigl|\rhoQM(x,t) - \Rbvec^T\Mab\eavec\bigr|^2\,dx
\]
\end{defn}

Projection loss is structural, not incidental.
The chain:
\[
  \Dista(\Ia)
  \;\geq\; \mathfrak{D}_C
  \;\geq\; \mathfrak{D}_R
  \;\geq\; \mathfrak{D}_M
\]
means that $\Lpsi > 0$ and $\Lrho > 0$ in general.
A model claiming $\Lpsi = 0$ exactly is either using all degrees of
freedom in $\Mab$ as free parameters, or it is wrong.


% ==================================================================
\section{Reconstruction Operator and Error}
\label{sec:recon}
% ==================================================================

\subsection{Reconstruction is not inversion}

Projection $\Pab : \Ha \to \Osp$ is generically non-invertible
because it is lossy.
Reconstruction produces an estimate, not a recovery of $\Ia$.

\begin{defn}[Reconstruction operator]
\label{def:Qba}
The \emph{reconstruction operator} is:
\[
  \boxed{\Qba : \Osp \to \Hahat}
\]
where $\Hahat$ is the reconstruction codomain (an approximation to
$\Ha$, not necessarily $\Ha$ itself).
The reconstructed identity estimate is:
\[
  \Ihat = \Qba(\Oba) = \Qba(\Pab(\Ia)).
\]
\end{defn}

\begin{defn}[Reconstruction error]
\label{def:reconerror}
Let $\dHa$ be a metric or pseudometric on $\Ha$.
The \emph{reconstruction error} is:
\[
  \boxed{
    \epsI = \dHa\!\bigl(\Ia,\;\Ihat\bigr)
    = \dHa\!\bigl(\Ia,\;\Qba\circ\Pab(\Ia)\bigr)
  }
\]
\end{defn}

Reconstruction is satisfactory when $\epsI < \delta$ for a
context-appropriate tolerance $\delta$:
\[
  \dHa\!\bigl(\Ia,\;\Qba\circ\Pab(\Ia)\bigr) < \delta.
\]

\begin{remark}
The phrase ``reconstruct the original'' is not used.
$\Qba$ produces an estimate $\Ihat$.
Whether $\Ihat \approx \Ia$ depends on how much of $\Ia$ was
preserved through $\Pab$, which is exactly the question that
$\etaab$ and $\epsI$ are defined to answer.
\end{remark}

\begin{remark}
The metric $\dHa$ on $\Ha$ is left unspecified here.
Specifying it requires structural assumptions about $\Ha$ that
the framework does not yet impose.
The definition is in place; the metric is a parameter to be
supplied by specific applications.
\end{remark}


% ==================================================================
\section{Compatibility as a Recoverability Gate}
\label{sec:compat}
% ==================================================================

\begin{defn}[Scale compatibility]
\label{def:compat}
\emph{Scale compatibility} $\chiab(\Ia)$ measures whether projection
from $\Sa$ to $\Sb$ is meaningfully recoverable for the specific
state $\Ia$.

\textit{Continuous form (recommended for IKK~III):}
\[
  \boxed{
    \chiab(\Ia)
    = \sigma\!\bigl(\alpha(\etaab(\Ia) - \etamin)\bigr),
    \qquad
    \sigma(z) = \frac{1}{1+e^{-z}}
  }
\]
where $\alpha > 0$ is the gate sharpness and $\etamin$ is the
minimum efficiency threshold.

\textit{Boolean limit:}
\[
  \chiab(\Ia)
  \;\to\; \mathbf{1}[\etaab(\Ia) > \etamin]
  \quad\text{as}\quad
  \alpha \to \infty.
\]
\end{defn}

\begin{remark}
The continuous form makes $\chiab$ differentiable in $\etaab$,
which is useful for gradient-based analysis.
The Boolean limit is the hard access gate: either the projection
is meaningful or it is not.
For IKK~III, the continuous form is preferred because it reflects
the gradual degradation of projection quality across scale
separation, rather than a sharp cutoff.
\end{remark}

\begin{remark}
Note that $\chiab$ is now a functional of $\Ia$ via $\etaab(\Ia)$,
not a scalar property of the scale pair $(\Sa,\Sb)$ alone.
Two identity states at the same source scale may be compatible
or incompatible with the same target scale depending on their
structure.
\end{remark}


% ==================================================================
\section{Wave-Particle Duality Reinterpreted}
\label{sec:duality}
% ==================================================================

\begin{prop}[Wave-particle duality as projection-domain switching]
\label{prop:duality}
Let $\Ia \in \Ha$ be a source identity state.
Let $\Pab^{\mathrm{particle}}$ denote projection into a domain
resolving localised identity features, and
$\Pab^{\mathrm{wave}}$ denote projection into a domain resolving
distributed amplitude and phase.
Then:
\[
  \Ia
  \;\xrightarrow{\;\mathcal{P}^{\mathrm{particle}}\;}
  O_{b,\mathrm{particle}}^{(a)},
  \qquad
  \Ia
  \;\xrightarrow{\;\mathcal{P}^{\mathrm{wave}}\;}
  O_{b,\mathrm{wave}}^{(a)}.
\]
Wave-particle duality is a consequence of applying different
projection maps to the same source identity, not a primitive
ontological property of the identity itself.
\end{prop}

\begin{center}
\fbox{\parbox{0.72\textwidth}{\centering\vspace{0.3cm}
  \[
    \text{wave-particle duality}
    \;=\;
    \text{projection-domain switching}
  \]
\vspace{0.3cm}}}
\end{center}

The standard quantum mechanical predictions survive because the
projected layer at $\Sb$ obeys the standard formalism in every
domain where that formalism has been validated.
IKK~III adds a proposed mechanism: $\II \to \Wab \to |\psi\rangle$.
Whether the mechanism is experimentally distinguishable from
standard quantum mechanics is the assignment of EVM~III.


% ==================================================================
\section{What IKK Must Still Earn}
\label{sec:debt}
% ==================================================================

\begin{epistemic}
Standard quantum mechanics remains the validated predictive layer.
IKK~III proposes a representational mechanism beneath it.
The framework earns legitimacy only if its typed projection mechanics
can recover known results and predict where recovery fails.
\end{epistemic}

\begin{evidenceitem}{1. Wavefunction reconstruction}
Does $\psirec(x,t) = \langle x\,|\,\Wab(\Ia(t))\rangle$
converge to $\psiQM(x,t)$ for benchmark systems?
Benchmark targets are listed in EVM~III.
$\Mab$ must be determined from the identity structure prior to
comparison.
\end{evidenceitem}

\begin{evidenceitem}{2. Density kernel fitting}
Does $\rhoQM(r,t) \approx \Rbvec^T\,\Mab\,\eavec$ with
$\Lrho \to 0$ as model resolution increases?
\end{evidenceitem}

\begin{evidenceitem}{3. Predictable error behavior}
Does $\epsI$ behave predictably as a function of $\Duab$ and
the structure of $\Ia$?
Mismatch that is reduced freely by parameter adjustment is
curve-fitting.
Mismatch that follows a predicted dependence is evidence.
\end{evidenceitem}

\begin{evidenceitem}{4. Scale dependence of efficiency}
Can $\etaab(\Ia) = f(\Duab, \Ia)$ be derived from the projection
structure and confirmed against simulation?
\end{evidenceitem}

\begin{evidenceitem}{5. Identified failure cases}
Where $\chiab(\Ia) \approx 0$ or $\etaab(\Ia) \ll 1$, the
framework must predict $\Lpsi \gg 0$ or $\Lrho \gg 0$.
A framework that cannot identify its own failure conditions is
not falsifiable.
\end{evidenceitem}

\begin{note}
Items~1--2 are addressed in EVM~I and EVM~III.
Items~3--5 are addressed in EVM~II and EVM~IV.
No claim of new physics is warranted until items~1 and~2 are
demonstrated.
\end{note}


% ==================================================================
\section*{Summary: Formal Projection Spine}
\addcontentsline{toc}{section}{Summary: Formal Projection Spine}
% ==================================================================

The following is the complete formal backbone of IKK~III.
Every symbol has a declared type and a named space.

\subsection*{Spaces}
\[
  \Ia \in \Ha,
  \qquad
  \Kab = \Cab(\Ha),
  \qquad
  \Hb \ni \psiab
\]

\subsection*{Typed maps}
\[
  \Cab : \Ha \to \Kab,
  \qquad
  \Rb : \Kab \to \Hb,
  \qquad
  \MXb : \Hb \to \ProbR
\]
\[
  \Wab = \Rb \circ \Cab : \Ha \to \Hb,
  \qquad
  \Qba : \Osp \to \Hahat
\]

\subsection*{Wave state and wavefunction}
\[
  |\psi_b^{(a)}\rangle = \Wab(\Ia),
  \qquad
  \psix = \langle x\,|\,\Wab(\Ia)\rangle
\]

\subsection*{Position measure and expectation}
\[
  \muX(\Omega)
  = \langle \Wab(\Ia),\;\EX(\Omega)\,\Wab(\Ia)\rangle_{\Hb}
  = \int_\Omega |\psix|^2\,d^3x
\]
\[
  \Xexp = \int_{\mathbb{R}^3} \mathbf{x}\;d\muX(\mathbf{x})
  = \langle\psi_b^{(a)}\,|\,\Xhat\,|\,\psi_b^{(a)}\rangle
\]

\subsection*{Density kernel}
\[
  \rhoab(r,t) = \Rbvec^T\,\Mab\,\eavec
\]

\subsection*{Efficiency, loss, and error}
\[
  \etaab(\Ia) = \frac{\Distb(\Pab(\Ia))}{\Dista(\Ia)},
  \qquad
  \Lpsi = \int |\psiQM - \psirec|^2\,dx,
  \qquad
  \Lrho = \int |\rhoQM - \Rbvec^T\Mab\eavec|^2\,dx
\]
\[
  \chiab(\Ia) = \sigma\!\bigl(\alpha(\etaab(\Ia)-\etamin)\bigr),
  \qquad
  \epsI = \dHa\!\bigl(\Ia,\;\Qba\circ\Pab(\Ia)\bigr)
\]

\subsection*{Central interpretive statements}
\begin{align*}
  \text{matrix mechanics}      &= \text{structured identity-state update} \\
  \text{wave mechanics}        &= \text{projected observable recovery} \\
  \text{wave-particle duality} &= \text{projection-domain switching}
\end{align*}

\bigskip

\begin{center}
\renewcommand{\arraystretch}{1.45}
\begin{tabular}{@{}lp{8.5cm}@{}}
  \toprule
  Symbol & Meaning \\
  \midrule
  $\Ha$          & Source identity space \\
  $\Hb$          & Target Hilbert space at $\Sb$ \\
  $\Kab$         & Coupled cross-scale information space \\
  $\Osp$         & Target observable space \\
  $\Itrue$       & Full primitive identity content \\
  $\Ia$          & Source identity state \\
  $\Ihat$        & Reconstructed identity estimate \\
  $\Cab$         & Cross-scale coupling map $\Ha\to\Kab$ \\
  $\Rb$          & Resolution map $\Kab\to\Hb$ \\
  $\Wab$         & Wave-state recovery map $\Ha\to\Hb$ \\
  $\MXb$         & Position measurement map $\Hb\to\ProbR$ \\
  $\Qba$         & Reconstruction operator $\Osp\to\Hahat$ \\
  $\EX$          & Spectral measure of $\Xhat$ on $\Hb$ \\
  $\muX$         & Target-scale position distribution of $\Ia$ \\
  $\Mab$         & Channel mixing matrix \\
  $\etaab(\Ia)$  & Projection efficiency functional \\
  $\chiab(\Ia)$  & Scale compatibility gate \\
  $\epsI$        & Reconstruction error $\dHa(\Ia,\Ihat)$ \\
  $\Lpsi$        & Wavefunction recovery loss \\
  $\Lrho$        & Density recovery loss \\
  $\uscl$        & Logarithmic scale coordinate $\ln(\ell/\ell_0)$ \\
  \bottomrule
\end{tabular}
\end{center}

\end{document}
